\documentclass[12pt,a4paper]{article}
\usepackage[UTF8,heading=true]{ctex}
\usepackage{geometry}
\usepackage{amsmath,amssymb}
\usepackage{booktabs,longtable,array,tabularx}
\usepackage{enumitem,hyperref,fancyhdr,xcolor,setspace}
\usepackage{ragged2e}

\newcolumntype{L}[1]{>{\RaggedRight\arraybackslash}p{#1}}
\ctexset{
    section = {
        name = {,、},
        number = \chinese{section},
        format = {\Large\bfseries\raggedright}
    }
}

\geometry{left=2.5cm,right=2.5cm,top=2.5cm,bottom=2.5cm}
\setlength{\headheight}{15pt}
\setstretch{1.5}
\setlength{\emergencystretch}{3em}
\hfuzz=20pt
\sloppy
\hypersetup{hidelinks}

\pagestyle{fancy}
\fancyhf{}
\fancyhead[C]{发明专利技术交底书}
\fancyfoot[C]{\thepage}

\begin{document}

\begin{center}
    {\Large\heiti 中国移动专利申请} \\[1em]
    {\Huge\heiti 技术交底书} \\[1em]
\end{center}

\noindent
{\setstretch{1.15}
\begin{tabular}{|L{2.8cm}|L{11.8cm}|}
\hline
公司编号 &  \\
\hline
发明名称 & 长距离量子中继器中的分布式纠缠纯化架构 \\
\hline
申报单位 & 研究院 \\
\hline
申报类型 & 发明 \\
\hline
发明人 & 待补充 \\
\hline
技术联系人 & 待补充 \\
\hline
注意事项 & 1．技术联系人应为深入了解本申请提案技术方案的技术人员。 \\
         & 2．请按照集团公司提供的本技术交底书模板逐项填写。 \\
\hline
\end{tabular}
}

\vspace{1cm}

\section{发明名称}

长距离量子中继器中的分布式纠缠纯化架构

\textit{英文名称：Distributed Entanglement Purification Architecture for Long-Distance Quantum Repeaters}

\section{技术领域}

本发明涉及量子中继器、量子互联网、纠缠分发、纠缠交换、量子存储与量子控制架构技术领域，尤其涉及一种在长距离量子通信链路中，面向\textbf{少量物理量子比特、有限量子存储相干时间和异步链路条件}而设计的分布式纠缠纯化架构及其控制方法。

本发明可适用于：
\begin{itemize}[noitemsep]
    \item 基于原子系综、NV 中心、离子阱、量子点或超导-光互连方案的量子中继节点；
    \item 长距离量子密钥分发、量子态分发、分布式量子计算与量子网络骨干链路；
    \item 具备光子接口、有限量子存储器以及经典前向/反馈链路的中继设备。
\end{itemize}

本发明不仅关注某个单独量子门，而是重点保护\textbf{中继节点内部资源组织方式、跨节点的纯化调度方式、纯化与交换的流水线耦合方式以及降低量子存储相干时间要求的控制策略}。

\section{现有技术的技术方案}

\subsection{量子中继器现有思路概述}

现有长距离量子通信通常无法采用经典放大器对量子态直接放大，因此需要量子中继器逐段建立纠缠，再通过纠缠交换将纠缠延伸到更长距离。当前公开技术的典型思路包括：
\begin{enumerate}[label=(\arabic*)]
    \item \textbf{基本量子中继链路}：先在相邻节点之间建立低保真度纠缠对，再通过多轮纯化提高保真度，最后执行纠缠交换；
    \item \textbf{多路复用量子中继}：通过频谱、时间或空间复用提高建立纠缠对的速率；
    \item \textbf{全光子量子中继}：尽量降低对长寿命量子存储的依赖；
    \item \textbf{大规模量子存储方案}：通过预置较多量子存储器以换取纯化成功概率。
\end{enumerate}

截至\textbf{2026年3月26日}，与本申请最接近的公开方案包括：
\begin{itemize}[noitemsep]
    \item \textbf{US7317574B2}：长距离量子通信相关方案，属于量子中继基础性专利；
    \item \textbf{US8135276B2}：量子中继器相关方案，强调中继节点结构；
    \item \textbf{US20250015901A1}：量子光网络中的量子中继器与方法；
    \item 学术上广泛引用的 \textbf{Briegel-D\"ur-Cirac-Zoller} 和 \textbf{D\"ur-Briegel} 系列量子中继与纠缠纯化工作。
\end{itemize}

\subsection{现有技术常见实现路径}

现有方案通常先批量建立多对低保真纠缠对，再选择两对或多对纠缠对进行双边门操作与测量，成功后保留高保真度纠缠对，失败则重新建立。该思路在理论上成立，但工程上往往依赖：
\begin{enumerate}[label=(\arabic*)]
    \item 较多的并行量子存储器；
    \item 节点间严格同步的控制节拍；
    \item 较长的量子存储相干时间；
    \item 重复失败时对大量历史纠缠资源进行整体丢弃。
\end{enumerate}

\section{现有技术的缺点及本申请提案要解决的技术问题}

现有技术至少存在以下问题：
\begin{enumerate}[label=(\arabic*)]
    \item \textbf{纯化资源开销大}。许多纯化架构默认每一级都需要同时持有多对纠缠对，导致对物理量子比特数量要求高。
    \item \textbf{相干时间要求苛刻}。如果等待足够多的纠缠对一起参与纯化，早期生成的纠缠对会因等待而退相干。
    \item \textbf{跨节点同步成本高}。多轮纯化与交换通常要求相邻节点严格同步，导致控制复杂度与时延上升。
    \item \textbf{失败恢复代价高}。一轮纯化失败可能导致整个链路段重新开始，资源浪费严重。
    \item \textbf{纯化与交换脱节}。很多方案把“纯化完成”与“交换执行”严格串行化，没有利用可重叠的时间窗口。
\end{enumerate}

因此，本申请要解决的核心问题是：
\begin{quote}
如何在长距离量子中继器中，以较少的物理量子比特和较短的量子存储相干时间，实现可分布式部署、可异步调度、可流水线执行的纠缠纯化与纠缠交换架构，从而快速得到高保真纠缠对。
\end{quote}

\section{本申请提案的技术方案的详细阐述}

\subsection{总体思路}

本发明提出一种\textbf{锚定对-牺牲对（anchor-sacrifice pair）}分布式纯化架构。每个中继节点不再要求同时缓存大量等价纠缠对，而是至少维护：
\begin{itemize}[noitemsep]
    \item 一个\textbf{锚定纠缠对}，作为当前准备被提升保真度的目标对；
    \item 一个或多个\textbf{候选牺牲纠缠对}，用于在合适窗口内参与纯化；
    \item 一个\textbf{局部辅助量子比特}，用于双边门与校验测量；
    \item 一个\textbf{经典控制与综合评分模块}，根据保真度估计、等待年龄和链路成功概率决定“纯化、交换或丢弃”。
\end{itemize}

为便于跨节点协同与工程实现，在一些实施例中，本发明对纠缠资源引入以下抽象标识：
\begin{itemize}[noitemsep]
    \item \textbf{链路段标识}：将长链路划分为相邻节点间的基本链路段，每段以 $(i,i+1)$ 标识；
    \item \textbf{层级标识（level）}：基础链路段建立的纠缠对为 $L0$，在中间节点执行一次纠缠交换后形成 $L1$，以此类推形成 $Lk$；
    \item \textbf{纠缠对标识（pair\_id）}：由链路段、层级、生成轮次、时间戳组成，用于压缩综合征消息关联与重放抑制；
    \item \textbf{年龄（age）}：用于描述纠缠对自“确认建链成功并写入存储量子比特”起的等待时间，用于刻画退相干风险。
\end{itemize}

本发明的架构核心是\textbf{事件驱动、局部协同、层级可递归}：任意节点只需与相邻节点交换少量压缩综合征与控制标志，即可在异步链路条件下滚动推进“建链-纯化-交换”，并在形成更高层级纠缠对后将其作为新的锚定对继续参与后续纯化/交换。

\subsection{节点资源组织方式}

在一个优选实施例中，每个中继节点至少包括：
\begin{enumerate}[label=(\arabic*)]
    \item 两个通信量子比特，用于与左右相邻节点建立纠缠；
    \item 两个存储量子比特，用于缓存锚定对和候选对；
    \item 一个辅助量子比特，用于执行局部门操作和校验；
    \item 一个光子接口或光子芯片子模块，用于纠缠对制备、贝尔态测量和经典结果上报；
    \item 一个本地控制器，用于计算评分、驱动门序列和管理纯化窗口。
\end{enumerate}

本发明的重要思想是：\textbf{用评分驱动的小批量滚动纯化替代大批量静态纯化}。

在资源复用方面，本发明优选采用“通信量子比特-存储量子比特”分离的使用方式：
\begin{enumerate}[label=(\arabic*)]
    \item 当一次建链尝试成功（例如基于光子探测的先验宣告 heralding 成功）后，将该纠缠态通过局部门操作交换（swap-in）写入存储量子比特，释放通信量子比特用于下一次建链尝试；
    \item 锚定对优先占用存储量子比特，候选牺牲对作为滚动缓冲占用另一存储量子比特；当牺牲对完成测量后，其占用立即释放，可继续写入下一对候选对；
    \item 辅助量子比特以“短时占用”为原则：仅在纯化门序列与校验测量窗口内占用，测量结束即释放用于下一轮纯化或交换辅助控制。
\end{enumerate}

\subsection{量子存储约束与资源复用策略}

量子存储器存在有限相干时间，且不同平台对退相干与操作误差的敏感性不同。本发明在控制层面优选显式建模以下约束，并据此做滚动调度：
\begin{enumerate}[label=(\arabic*)]
    \item \textbf{存储超时约束}：当锚定对年龄 $\Delta t_a$ 超过阈值 $T_{\mathrm{drop}}$（例如与 $T_2$ 同阶）时，优先触发局部丢弃与重建，避免“等待更久但最终不可用”的资源浪费；
    \item \textbf{操作窗口约束}：纯化与交换均需要双边门、测量与经典结果对齐。本发明优选将其抽象为可重叠的时间窗口：窗口内锁定相关量子比特并执行门序列，窗口外释放资源并继续建链；
    \item \textbf{局部锁定与租约（lease）}：为避免异步环境下两端节点对同一纠缠对做出冲突操作，本发明可在压缩综合征消息中携带短租约字段（例如“本轮纯化会话 id + 有效时间”），租约到期自动释放，降低死锁风险；
    \item \textbf{预测保真度约束}：对“可交换门限”不以当前估计保真度为唯一依据，而优选以“预计在完成交换门序列与经典往返后仍满足门限”的预测保真度为依据，从而将经典往返时延与退相干共同纳入决策。
\end{enumerate}

\subsection{分布式评分与决策机制}

设某一锚定纠缠对的估计保真度为 $F_a$，年龄为 $\Delta t_a$，量子存储器相干时间常数为 $T_2$，候选牺牲对的估计保真度为 $F_c$，当前链路重建代价估计为 $J_r$。系统定义决策评分：
\[
Q = w_1 F_a + w_2 F_c - w_3 \frac{\Delta t_a}{T_2} - w_4 J_r - w_5 P_{\mathrm{fail}},
\]
其中 $P_{\mathrm{fail}}$ 表示根据历史门保真度、探测成功率和链路状态推断的本轮纯化失败概率。

在一些实施例中，上述 $F_a$、$F_c$ 不是单次测量的瞬时值，而是由节点基于历史校验测量结果、器件标定参数和退相干模型得到的区间估计；$P_{\mathrm{fail}}$ 可结合本轮门序列长度、双边门保真度、探测效率与暗计数率等因素进行预测。评分权重 $w_i$ 与门限 $\theta_p,\theta_s,\theta_d$ 可由离线仿真优化得到，也可在在线运行中基于吞吐与失败率进行自适应更新。

若满足：
\[
Q \ge \theta_p,
\]
则触发纯化；若锚定对保真度已达到交换门限 $\theta_s$，则直接进入纠缠交换；若 $Q < \theta_d$，则丢弃年龄过大或价值过低的锚定对并重建。

上述机制使系统不必等待“所有节点均完全就绪”才开始纯化，而是允许相邻链路在可容忍的异步状态下局部推进。

\subsection{纯化与交换的流水线执行}

本发明的一个核心创新在于：将纯化与交换由传统的串行阶段式处理，改为\textbf{窗口重叠式流水线}处理：
\begin{enumerate}[label=\textbf{S\arabic*.}]
    \item 相邻节点之间并行产生初始纠缠对，并以最早成功的一对作为锚定对；
    \item 后续新产生的纠缠对不立即平等缓存，而是按评分规则标记为候选牺牲对；
    \item 一旦锚定对尚未达到交换门限且存在高评分候选对，则在当前节点和相邻节点执行双边门序列与测量；
    \item 纯化成功后，输出的新锚定对更新保真度和时间戳，并立即评估是否可以发起更高层的纠缠交换；
    \item 纯化失败时，仅丢弃受影响的局部对，不要求整条长链路回退到初始状态；
    \item 当中间节点的左右两侧锚定对均达到交换门限时，立即执行贝尔态测量完成交换，而无需等待远端节点完成下一轮纯化；
    \item 交换结果和纯化结果以压缩综合征形式沿经典链路广播，用于上层节点更新评分。
\end{enumerate}

为进一步明确时序关系，在一个优选实施例中，节点的控制流程可细化为三条并行循环（同一组物理量子比特在不同窗口内被复用）：
\begin{enumerate}[label=(\arabic*)]
    \item \textbf{建链循环}：通信量子比特持续触发与相邻节点的纠缠对生成尝试；一旦成功即将纠缠态写入存储量子比特并生成 pair\_id，随后立刻释放通信量子比特进入下一次尝试；
    \item \textbf{纯化循环}：当评分模块判定可纯化时，向相邻节点发送 \texttt{PurifyStart(pair\_id\_a, pair\_id\_c, session\_id, lease)}，两端在租约内对齐门序列并执行双边门与测量；仅交换必要测量结果即可决定保留/丢弃，并更新锚定对的帕里框架（Pauli frame）与保真度估计；
    \item \textbf{交换循环}：当某中间节点同时持有“左锚定对”和“右锚定对”且满足交换门限时，在本节点执行贝尔态测量（BSM）完成纠缠交换，并将交换结果以 \texttt{SwapResult(session\_id, bsm\_bits, timestamp)} 形式通知两侧端点节点，端点节点据此更新长距离纠缠对的帕里框架与状态标记。
\end{enumerate}

上述并行循环使得系统可以在等待经典确认或等待更高层锚定对就绪的间隙继续“生成候选牺牲对”，从而提高吞吐；同时，通过将失败恢复限制在“当前链路段或当前会话”，避免上层长链路整体回退造成的资源雪崩式浪费。

\subsection{压缩综合征与分布式协调}

传统方案常把大量历史门结果和原始测量记录向上汇总，本发明则优选采用\textbf{压缩综合征消息}，其至少包括：
\begin{itemize}[noitemsep]
    \item 锚定对标识；
    \item 最近一轮纯化/交换的成功失败标志；
    \item 必要的帕里蒂校验测量结果；
    \item 时间戳和年龄信息；
    \item 更新后的保真度区间估计。
\end{itemize}

在一些实施例中，压缩综合征消息还可包括：层级标识 level、会话标识 session\_id、租约字段 lease、以及用于重放抑制的单调递增计数器。对于纠缠交换，若采用标准 BSM，其结果通常仅需 2 比特即可更新 Pauli frame；对于纯化校验，仅需上报决定“保留/丢弃”与必要的帕里蒂结果即可驱动上层评分更新。通过这一方式，可以将跨节点控制信息压缩为“决策充分信息”，而不是传输全部原始测量轨迹。

通过仅交换压缩后的决策充分信息，可以降低经典控制开销，并便于在 FPGA/ASIC 或光子芯片控制器中实现。

\subsection{量子门序列的优选方式}

本发明不限定具体平台。在一个优选实施例中，纯化门序列可为局部 CNOT、受控相位门、单比特旋转和条件测量的组合；在另一实施例中，可由集成光子芯片中的干涉网络、可调相移器和单光子探测器完成等效操作。

为降低实现不确定性，以下给出若干\textbf{可选且可替换}的门序列示例（不构成对保护范围的限制）：
\begin{enumerate}[label=(\arabic*)]
    \item \textbf{双边纯化序列示例}：两端节点分别对“锚定对 + 牺牲对”执行双边 CNOT（或等效受控门），随后测量牺牲对并交换测量比特；当帕里蒂条件满足时保留锚定对并完成一次纯化，否则丢弃本轮涉及的局部资源并回到建链循环。
    \item \textbf{交换（BSM）序列示例（物质量子比特实现）}：中间节点对其持有的两端量子比特执行 CNOT 与 Hadamard（或等效基变换），随后对两比特测量得到 2 比特 BSM 结果，并沿经典链路发送给端点节点用于 Pauli frame 更新。
    \item \textbf{交换（BSM）序列示例（光子芯片实现）}：中间节点将两侧到达的光子态注入干涉网络，通过可调相移器与分束器阵列实现贝尔态测量的等效投影，并用单光子探测器输出 2 比特结果；该结果同样可用于端点节点的 Pauli frame 更新。
    \item \textbf{片上控制实现示例}：评分模块与状态机可在 FPGA/ASIC 内实现事件驱动调度；与集成光子芯片协同时，可将 \texttt{PurifyStart}、\texttt{SwapReady} 等消息直接映射为片上时钟域内的触发信号，以缩短纯化窗口与交换窗口的控制开销。
\end{enumerate}

本发明保护的重点不在于某个唯一门序列，而在于：
\begin{enumerate}[label=(\arabic*)]
    \item \textbf{根据评分决定何时纯化、何时交换、何时丢弃}；
    \item \textbf{锚定对与候选牺牲对的资源组织结构}；
    \item \textbf{纯化窗口与交换窗口的重叠流水线机制}；
    \item \textbf{压缩综合征驱动的分布式协调方式}。
\end{enumerate}

\subsection{优选实施例}

设一条 1000 公里量子中继链路被分为若干 50 公里到 100 公里的基本链路段。每个中继节点只配置 5 个左右的关键物理量子比特。系统首先在每一基本链路段上生成初始纠缠对：
\begin{enumerate}[label=(\arabic*)]
    \item 最先成功的一对被登记为锚定对；
    \item 随后生成的新对作为候选牺牲对；
    \item 节点根据锚定对老化速度与候选对质量决定是否立即纯化；
    \item 某一中间节点左右两侧都达到交换门限后，立即执行交换；
    \item 失败恢复仅限于局部链路段，不要求全局回退。
\end{enumerate}

在该实施例中，系统可以在有限存储条件下持续滚动输出更高保真度的长距离纠缠对。

\subsection{本发明的有益效果}

与现有技术相比，本发明至少具有以下效果：
\begin{enumerate}[label=(\arabic*)]
    \item \textbf{降低物理量子比特需求}：通过锚定对-牺牲对组织方式，用较少缓存资源实现连续纯化；
    \item \textbf{降低相干时间压力}：通过滚动式、小批量、评分驱动的纯化，减少等待多个纠缠对齐套可用的时间；
    \item \textbf{增强异步容忍度}：无需全局严格同步，可按局部链路状态推进；
    \item \textbf{减少失败回退范围}：局部失败局部恢复，提升总体吞吐；
    \item \textbf{便于工程实现}：适合与 FPGA、ASIC 或集成光子芯片的控制逻辑配合落地。
\end{enumerate}

\section{本申请建议重点保护的内容}

\subsection{节点状态机与触发逻辑}

建议在正式申请中把节点控制状态机明确描述为至少包括以下状态：
\begin{enumerate}[label=(\arabic*)]
    \item \textbf{等待建链状态}：持续尝试生成新的基础纠缠对；
    \item \textbf{锚定保持状态}：当首个合格纠缠对建立后，将其登记为锚定对并持续监测其年龄与保真度；
    \item \textbf{局部纯化状态}：当存在高评分候选牺牲对且满足纯化门限时，执行纯化门序列；
    \item \textbf{交换就绪状态}：当左右两侧锚定对均达到交换门限时，等待并触发贝尔态测量；
    \item \textbf{局部恢复状态}：当局部纯化或交换失败时，仅回退受影响的链路段而不全局回退。
\end{enumerate}

对应地，建议明确下列触发条件：
\begin{enumerate}[label=(\arabic*)]
    \item 当 $Q \ge \theta_p$ 且锚定对年龄未超限时，触发纯化；
    \item 当 $F_a \ge \theta_s$ 且相邻两侧均满足交换条件时，触发交换；
    \item 当 $\Delta t_a > T_{\mathrm{drop}}$ 或 $Q < \theta_d$ 时，触发局部丢弃与重建；
    \item 当经典控制信道报告上游已形成更高层候选链路时，可提高本地交换优先级。
\end{enumerate}

建议重点保护以下技术点：
\begin{enumerate}[label=(\arabic*)]
    \item 锚定对-牺牲对式的中继节点量子资源组织结构；
    \item 基于保真度、年龄、失败概率和回退代价的纯化/交换评分决策方法；
    \item 纯化与纠缠交换窗口重叠的流水线控制方法；
    \item 压缩综合征驱动的分布式协调与状态更新方法；
    \item 在少量物理量子比特条件下的中继节点实现架构。
\end{enumerate}

建议正式申请时按照“方法 + 节点装置 + 网络系统 + 控制程序”四层布局：
\begin{enumerate}[label=(\arabic*)]
    \item \textbf{方法权利要求}：限定评分计算、状态切换、纯化执行和交换执行流程；
    \item \textbf{节点装置权利要求}：限定锚定对存储单元、候选对管理单元、辅助量子比特和本地控制器之间的结构关系；
    \item \textbf{网络系统权利要求}：限定多个中继节点通过压缩综合征消息形成分布式协调；
    \item \textbf{程序或介质权利要求}：覆盖评分驱动调度逻辑的软件实现。
\end{enumerate}

\section{附图建议}

建议在正式申请中至少配置以下附图：
\begin{enumerate}[label=(\arabic*)]
    \item 量子中继节点硬件结构图：展示通信量子比特、存储量子比特、辅助量子比特和光子接口；
    \item 锚定对-牺牲对资源组织示意图：突出本发明与传统批处理纯化的差异；
    \item 节点状态机图：展示等待建链、局部纯化、交换就绪和局部恢复状态切换；
    \item 纯化与交换窗口重叠时序图：展示流水线执行机制；
    \item 压缩综合征消息交互图：展示分布式协调方式。
\end{enumerate}

\section{可替代实施方式}

本发明可以运行于离子阱、NV 中心、量子点、自旋光子接口或其他具备量子存储与光学链路的硬件平台；评分函数既可由解析模型给出，也可由机器学习预测器给出；压缩综合征可采用固定长度比特字段，也可采用带置信区间的结构化消息。

\section{结论}

本发明针对量子中继器难以工程化的核心瓶颈，提出了分布式纠缠纯化架构。其本质创新不在于抽象的“再做一次纯化”，而在于\textbf{如何在有限资源、有限相干时间与分布式异步条件下，持续、滚动、局部化地完成纯化与交换}。该方案有望成为未来量子互联网骨干网络的重要底层控制专利。

\end{document}
